
% ════════════════════════════════════════════
% 具体分析（5.1.1.具体分析.tex，三级子块）
% ════════════════════════════════════════════
\subsubsection{问题一的分析与建模思路}

问题一属确定性线性优化问题，采用线性规划对微网一天的购电与储能充放电策略联合优化。

先将附件一给出的单日分时电价、小区负载与光伏预测数据读取并时间对齐，整理为一天 144 个 10 分钟时间间隔上的功率序列；再依据储能设备的容量、充放电功率与效率约束，把一天内的能量平衡关系与储能荷电状态递推关系写成线性约束。模型以全天购电费用最小为目标函数，加入储能设备在 0:00 与 24:00 储电量相等的日循环约束，并按附录1给定值固定 0:00 储电量为 6000 kWh，构成标准线性规划；调用求解器求得每个时间间隔的购电功率与储能充放电功率，进而得到全天购电量、购电费以及各时段储能充放电量。
